
% SGA 3，第 VII 篇_B；英文重建本第一轮的中文译本。
% 范围： 节 3.2 从开头至命题 3.2,
% 依据文献局部页 67--68, 止于下节之前： 3.2.1.
% 合并阅读本页： 580--581. 重版印刷页： 569--570.
% 范围内原版边码： 532--534.
% 范围内编者注： 119--123.
% 依据： Polo--Gille Exp7B-13oct24.pdf,
% SHA-256 D303690E0195654BBE7C3DAA5D89EA23E24CD709A3575217234AC19AE8836694.
% Jacob Reinhold 的 CC BY 4.0 Markdown 仅用作对照来源，不作为依据。
% 本组件无图；所有数学内容均为原生 LaTeX。
\subsection*{3.2}
\phantomsection
\label{sga3:VIIB:3.2}
现在设 \(C\) 是 Artin 环，其根为 \(\mathfrak r\)。于是对每个\textup{（结合且含幺的）}\(C\)-代数 \(U\)，理想 \(\mathfrak rU\) 都由幂零元组成。若 \(x\in\mathfrak rU\)，置
\[
  \exp_U(x)
  =
  1+x+\frac{x^2}{2!}+\cdots.
\]
\footnote{\label{sga3:VIIB:nde:119}若 \(x,x'\in\mathfrak rU\) 交换，则
\[
  \exp_U(x+x')=\exp_U(x)\exp_U(x').
\]
}
这给出从 \(\mathfrak rU\) 到 \(1+\mathfrak rU\) 的双射；其逆将 \(1+y\in1+\mathfrak rU\) 映为
\[
  \log_U(1+y)
  =
  y-\frac{y^2}{2}+\frac{y^3}{3}-\cdots.
\]
此外，映射 \(\exp_U\) 显然关于 \(U\) 是函子性的。\footnote{\label{sga3:VIIB:nde:120}即对每个 \(C\)-代数同态 \(\phi:U\to V\)，
\[
  \phi\bigl(\exp_U(x)\bigr)=\exp_V\bigl(\phi(x)\bigr).
\]
}
仍设 \(C\) 为 Artin 环，进一步设 \(U\) 赋有 \(C\) 上的双代数结构\textup{（参见\ \SGARefLink{sga3:VIIB:2.2}{2.2}）}。则对每个本原元素 \(x\in\mathfrak rU\)\textup{（参见 VII\textsubscript{A}，\SGARefLink{sga3:VIIA:3.2.3}{3.2.3}）}，有
\[
\begin{aligned}
  \Delta_U\bigl(\exp_U(x)\bigr)
  &=
  \exp_{U\otimes_CU}\bigl(\Delta_U(x)\bigr)
  \\
  &=
  \exp_{U\otimes_CU}(x\otimes1+1\otimes x)
  \\
  &=
  \exp_{U\otimes_CU}(x\otimes1)\,
  \exp_{U\otimes_CU}(1\otimes x)
  \\
  &=
  \bigl(\exp_U(x)\otimes1\bigr)
  \bigl(1\otimes\exp_U(x)\bigr)
  \\
  &=
  \exp_U(x)\otimes\exp_U(x).
\end{aligned}
\]
因此，双射 \(\exp_U\) 将 \(\mathfrak rU\) 的本原元素映为满足
\[
  \Delta_U(z)=z\otimes z.
\]
的元素 \(z\in1+\mathfrak rU\)。反之，若 \(z\) 满足这些条件，则置 \(x=\log_U(z)\)，上述计算表明
\[
  \exp_{U\otimes_CU}(x\otimes1+1\otimes x)
  =
  z\otimes z
  =
  \Delta_U\bigl(\exp_U(x)\bigr)
  =
  \exp_{U\otimes_CU}\bigl(\Delta_U(x)\bigr);
\]
故
\[
  x\otimes1+1\otimes x=\Delta_U(x).
\]
\footnote{\label{sga3:VIIB:nde:121}编者补充了上述论证的细节，并添加了下句。}
进一步，若 \(z\in1+\mathfrak rU\) 满足 \(\Delta_U(z)=z\otimes z\)，则
\[
  \varepsilon_U(z)^2=\varepsilon_U(z).
\]
由于 \(\varepsilon_U(z)\) 可逆\textup{（因为 \(z\) 可逆，且 \(\mathfrak r\) 幂零）}，可得 \(\varepsilon_U(z)=1\)。
特别地，设 \(L\) 是 \(C\) 上的平坦 Lie 代数，取 \(U\) 为 \(L\) 在 \(C\) 上的包络代数 \(U(L)\)，并将 \(L\) 与其在 \(U\) 中的典范像等同。由\SGARefLink{sga3:VIIB:3.1:lemma}{引理 3.1}，\(L\) 是 \(U\) 的本原元素集合：事实上，在 \(C\) 的每个局部分量上，\(L\) 都是自由模。考察 \(C\)-形式群
\[
  \mathscr G(L)=\operatorname{Spf}^{*}U(L),
\]
其协变双代数为 \(U=U(L)\)\textup{（参见\ \SGARefLink{sga3:VIIB:2.6.2:prop}{2.6.2}）}。设 \(\mathfrak m\) 是 \(C\) 的一个极大理想，置
\[
  \overline C=C/\mathfrak m,
  \qquad
  \overline U=U\otimes_C\overline C.
\]
由于 \(\mathscr G(L)\) 无穷小\textup{（同上）}，\(\overline U\) 的单位元是唯一满足下式的元素 \(\overline z\in\overline U\)：\footnote{\label{sga3:VIIB:nde:122}编者补充了原文遗漏的条件 \(\varepsilon_{\overline U}(\overline z)=1\)。}
\[
  \varepsilon_{\overline U}(\overline z)=1,
  \qquad
  \Delta_{\overline U}(\overline z)
  =
  \overline z\otimes\overline z.
\]
由此可得，每个满足
\[
  \varepsilon_U(z)=1,
  \qquad
  \Delta_U(z)=z\otimes z
\]
的 \(z\in U\) 必属于 \(1+\mathfrak rU\)。
最后，由于
\[
  L\cap\mathfrak rU
  =
  \mathfrak rL
  =
  L\otimes_C\mathfrak r,
\]
可知 \(\exp_U\) 定义了从 \(L\otimes_C\mathfrak r\) 到群 \(\mathscr G(L)(C)\) 的双射。将此归纳如下。
\medskip
\noindent\textbf{命题 3.2。}
\phantomsection
\label{sga3:VIIB:3.2:prop}
{\itshape 设 \(k\) 是包含 \(\mathbb Q\) 的伪紧环，\(\mathbf L\) 是 \(\mathbf O_k\) 上的平坦 Lie 代数。
\begin{enumerate}
\item[\textup{(i)}]
对 \(\operatorname{Alf}/k\) 的每个对象 \(C\)，以 \(\mathfrak r(C)\) 表示其根。则映射
\[
  \exp_{U(\mathbf L(C))}:
  \mathbf L(C)\otimes_C\mathfrak r(C)
  \longrightarrow
  \mathscr G(\mathbf L)(C)
\]
是双射，且关于 \(C\) 和 \(\mathbf L\) 是函子性的。
\item[\textup{\SGAInlineRefTarget{sga3:VIIB:3.2:prop:item:ii}(ii)}]
自然态射
\[
  \mathbf L
  \longrightarrow
  \operatorname{Lie}\bigl(\mathscr G(\mathbf L)\bigr)
\]
\textup{（参见\ \SGARefLink{sga3:VIIB:2.6.2:prop}{2.6.2}）}是 \(\mathbf O_k\) 上 Lie 代数的同构。
\end{enumerate}}
\footnote{\label{sga3:VIIB:nde:123}第\textup{\SGARefLink{sga3:VIIB:3.2:prop:item:ii}{(ii)}} 项由编者添加，它直接来自 \SGARefLink{sga3:VIIB:3.1:lemma}{3.1}。}
